In this work, we formalize and study the question: can we train models that learn different classes of functions in-context? We show that Transformer models trained from scratch can in-context learn the class of linear functions, with performance comparable to the optimal least squares estimator, even under distribution shifts. Moreover, we  show that in-context learning is also possible for sparse linear functions, decision trees, and two-layer neural networks; learning problems which are solved in practice with involved iterative algorithms such as gradient descent. %Moreover, the model trained to in-context learn two-layer neural networks is also able to in-context learn linear functions (for which it is not explicitly trained). 

At the same time, understanding the implications of our results for language models requires further investigation. A pertinent question regarding the in-context learning capabilities of language models is how they leverage in-context examples \citep{min2022rethinking}. Our results demonstrate that Transformers can encode complex learning algorithms that utilize in-context examples in a far-from-trivial manner. In fact, this is the case for standard Transformer architectures trained with standard optimization procedures. The extent to which such non-trivial in-context learning behavior exists in large language models is still open, but we believe that our work takes a step towards formalizing and investigating this question.

Our work lays the groundwork for several future directions.

\paragraph{Complexity of in-context learning.} We empirically show that model capacity helps in performing in-context learning accurately and robustly. This raises the question: How does the in-context learning loss \eqref{eq:in-context_loss} depend on the complexity of the function class $\mathcal{F}$, the capacity of model $M$, and the number of prompts used to train $M$. Even the right notion of complexity of $\mathcal{F}$ is unclear and may depend on the model family. Understanding this question for models explicitly trained to perform in-context learning may suggest an upper bound for the in-context learning performance of models such as GPT-3 that have not been explicitly trained for this purpose.

\paragraph{Curriculum learning.} Within our framework, there is natural notion of curriculum learning where during training, we gradually increase the complexity of the function class learned in-context. This leads to drastic speed-ups in training. What is the reason behind such a speedup? Are similar speedups also possible for training large language models? Understanding these questions can have implications for training of models on large real-world datasets, potentially reducing the time and energy used for training.

\paragraph{Inductive bias of model families.} Our framework presents an opportunity to understand and compare the inductive biases of different model families (e.g., Transformers vs.\ LSTMs) in a well-defined setting. For instance, a concrete question is: Are there function classes that are easier to in-context learn using Transformers but harder for LSTMs and vice-versa?

\paragraph{Understanding the learning algorithms encoded in Transformers.}
The models we train are able to perform in-context learning, and are thus themselves encoding learning algorithms. A worthwhile research
direction would be to investigate the internal workings of these models and better understand the exact
learning algorithms that they encode.
Moreover, for settings such as decision trees, we do not have a good understanding of what the optimal learning algorithms are.
Nevertheless, in Section~\ref{sec:beyond_linear} we found that Transformers are able to discover sample efficient algorithms when being trained to perform in-context learning.
This suggests an intriguing possibility where we might be able to reverse engineer the Transformer to obtain better learning algorithms for such problems.